\scanpage{59}
% archive sheet number: 58
Ainsi $\unsure{\psi|\varphi}$ \uline{est un 1-cocycle}, quand on fait opérer $\operatorname{Gal}(k/k_0)$ sur $G(k)$ de la façon évidente (par conjugaison). Il correspond à un fibré principal homogène $P_{v,k}$ sous $G(k)$, avec opérations de $\operatorname{Gal}(k/k_0)$ sur $(P_v,G(k))$ \deleted{\unsure{compatible avec le gpe}}. Il définit \inserted{par descente} un fibré principal homogène \inserted{$P_{v,k_0}$} sous $G_{k_0}$, avec un \uline{\unsure{splittage|splitting|splitt.} donné} sur $k$.

\margin{\uline{À vérifier} OK.}
% slanted
On peut d'ailleurs dire ce que c'est que $P$ : c'est le fibré principal hom des isoms \inserted{\unsure{du|de} foncteur} $\xi_{M,\deleted{\unsure{k}}\infty}$ avec $\xi_{M,v}^{\mathbb Z}\otimes_{\mathbb Z} k_0$, dont une trivialisation est en effet obtenue \ldots\ldots

Par la conjecture \inserted{de transcendance} N\textsuperscript{o} 1, \deleted{\unsure{cette} \unsure{triv.|fini} \unsure{l'élém}} \unsure{ce|le} point ainsi donné sur $P_{v\,k}$ est \uline{générique}, revient à dire que l'application
\margin{$\pi$ transcendant sur $k_0$\ldots.}
% slanted
\diagram{\begin{tikzcd}[ampersand replacement=\&]
\operatorname{Gal}(k/k_0)\arrow[r,"{\varphi_v}"]\&G(k)\,[\subset\mathrm{Gl}(M_k(\infty))]
\end{tikzcd}}
\deleted{est} a une image Zariski-dense. [N.B. cela peut s'écrire,
